\documentclass{zhvt-classic}

% 使用本地源流明体（GenRyuMin2 TW）按文件路径加载，不依赖系统字体名
% 版框（内框／界栏／书口侧栏／鱼尾）改用与版心印章同一印泥棕色：
%   印章图片 yshs.png 实色 = RGB(135,68,52) = HTML 874434
\usepackage{xcolor}
\usepackage{tikz}
% sealbrown 双注册：xcolor 供 TikZ/pgf，L3 供 cls 版框
\definecolor{sealbrown}{HTML}{874434}
\ExplSyntaxOn
\color_set:nnn { sealbrown } { HTML } { 874434 }
\ExplSyntaxOff

\zhvtset{
  grid color       = sealbrown,
  main font        = GenRyuMin2TW-SB,
  main font path   = C:/Users/Administrator/texmf/fonts/opentype/public/GenRyuMin2TW/,
jiazhu font      = GenRyuMin2TW-M,
  jiazhu font path = C:/Users/Administrator/texmf/fonts/opentype/public/GenRyuMin2TW/,
  % 页码框容纳 4 字：正常页「九十七」+前置空 1 格=4 行恰好；
  % 100+ 页「一百零一」共 5 行超出，末字「一/二/三」被框裁掉（2026-10-04 实测），故放宽到 5
  shukou page number chars = 5,
  % 书口书名框高从默认 6 字减到 4 字，把省出的 2 字高拨给 chapter 框，
  % 使 6 字书名+卷名放进 9 字高的 chapter 框后不再压线（2026-10-05）
  shukou title chars = 4,
}

\title{史记}[卫康叔世家]
\author{汉·司马迁}
\maker{闲叟漫记·重製}


% ===== 索隐述赞样式（2026-11-08 四次定稿 · 愚大人亲手改定参数）=====
% ⚠ 竖排前提：自定义字体家族必须带 RawFeature={vertical:+vert}；
%   simhei.ttf 加 +vhal 会整体加载失败静默回退，只写 +vert；simfang 同理。
\newCJKfontfamily{\shuzanfont}[Path=C:/Windows/Fonts/,RawFeature={vertical:+vert}]{simhei.ttf}
\newCJKfontfamily{\shuzanbody}[Path=C:/Windows/Fonts/,RawFeature={vertical:+vert}]{simfang.ttf}
% 题头：152x30pt 圆角 6pt 棕底白字；baseline (0,2) 使整框落入界行格子不压版框线；
% 白字节点 x 不补偿、y +8.2pt（像素级墨迹实测 dx/dy≈0，L/R 留白约 20pt）。
\newcommand{\shuzanhead}[1]{%
  \begin{tikzpicture}[baseline={(0,2)}]
    \fill[sealbrown, rounded corners=6pt]
      (-76pt,-15pt) rectangle (76pt,15pt);
    \node[text=white, inner sep=0pt, anchor=center] at (0pt,8.2pt)
      {\shuzanfont\fontsize{23}{23}\selectfont #1};
  \end{tikzpicture}%
}
% ★ 修正书口卷名位置：chapter 框起始位置改回 cls 默认值，避免全局影响。
%   只在每个卷首页通过 \juan 命令里的空命令微调该页位置（安全，不会跨页污染）。
\ExplSyntaxOn
\AtBeginDocument{
  \dim_set:Nn \l__zhvt_shukou_chapter_title_start_dim
    { \l__zhvt_shukou_title_height_dim + \l__zhvt_font_size_dim * 2 }
  \dim_set:Nn \l__zhvt_shukou_page_number_start_dim
    { \l__zhvt_shukou_chapter_title_start_dim + \l__zhvt_shukou_chapter_title_height_dim }
}
\ExplSyntaxOff

\begin{document}

% ★★ 版心（书口）排布覆写：把书名放到鱼尾之下（2026-10-04）
%
%   cls 默认版心自上而下：书名（title 静态框，在鱼尾之上）→ 鱼尾 → 篇名（chapter 动态框）→ 页码
%   本模板改为：          鱼尾（位置不动）→ 书名 → 篇名 → 页码分割线 → 页码
%
%   手法：把 cls 放在 title 静态框里的书名，改放进 chapter 动态框（鱼尾之下）；
%         title 静态框留空。书名取 \l__zhvt_title_str，即 \title{...} 的第一参数，
%         所以换书时只改 \title，不必在此重复写书名。
%
%   ⚠ 三条必须遵守的写法（2026-10-04 实测）：
%     ① 汉字必须用 \hbox{...} 包住。flowfram 会把框内容装进一个宽 14.4mm（= 一个 line_sep
%        = 书口栏宽）的纵向 minipage 并按该宽度折行，而整页靠 PDF /Rotate 90 显示、
%        框内文字本身不变向 —— 不包 \hbox 时「围炉夜话」4 字会被逐字折成 4 行，
%        在版面上就成了横排 4 字并横向溢出版框。
%     ② 竖排下**一个 \hbox 就是一整列**。书名与篇名要同处一列，必须放进同一个 \hbox；
%        分成两个 \hbox 时，后一个会被当成新的一列、甩到隔壁栏去
%        （实测「序」跑到书口栏外侧 14.4mm 的正文列上）。中间的 \quad 在框内是横向空白，
%        显示出来正好是书名与篇名之间的纵向间隔。
%     ③ ★ 内容盒要写成 **与框等宽的定宽盒** `\hbox to \l__zhvt_shukou_chapter_title_height_dim{…\hfil}`。
%        实测：chapter 框是把内容盒沿纵向**居中**摆放的，内容宽度一变，书名就跟着漂：
%        自然宽度的 \hbox（书名 38.4mm）落在 x=157.2mm，加篇名后（57.6mm）落在 x=152.4mm，
%        目录页（无篇名）与正文页的书名就差半格。做成等宽盒后，无论有没有篇名，
%        内容固定从框顶起排，实测两页均落在 x=142.8mm，不再漂。
\ExplSyntaxOn
% 把书口 chapter 框（书名+篇名）起始位置下移 1 字高，避免 4 字卷名顶到鱼尾
% ⚠ cls 的 \__zhvt_calculate_dimensions: 在 \AtBeginDocument 里执行，会重算该 dim；
%   必须把覆写挂到 begin document 之后，否则会被二次覆盖（2026-10-05 实测）。
% (旧方案已移到导言区，直接修改 chapter 框起始位置)
% 卷标题命令 #1=书口篇名(鱼尾) #2=正文卷名（含卷序，卷名内空格用 ~）
%
% ★★ 目录页两项修订（2026-10-05，愚大人指出）★★
%   ① 页码起点统一对齐：
%      原状 —— 目录项写成 \par\hspace{4em}卷名\hspace{4em}页码，页码起点 = 卷名实宽 + 4em。
%      卷名实宽随字数变：「卷一~梁惠王上」6 字 +~、「卷七~离娄上」5 字 +~，
%      差 1 个字位 → 卷七~卷十 4 条的页码整体比其余 10 条高 27.21pt（一个字），看着上下不齐。
%      改法 —— 卷名补全角空格（宽 = 1 字 = 27.215pt）补足到 **7 字符位**再写进目录，
%      14 条页码起点遂全部落在同一条竖线上（起点对齐；末字随页码字数自然参差）。
%      ⚠ 只对目录里的副本 \l_tmpc_tl 补，正文卷名仍用 \l_tmpa_tl 原样 —— 正文像素不受影响。
%   ② 卷名整体上移半字位：前导 \hspace 由 4em 改 3.5em。
%      实测原状目录块在版框内偏下（天头留白 139.6pt / 地脚留白 110.7pt），
%      上移 0.5em = 13.6pt 后上下基本均分，视觉居中。
\NewDocumentCommand { \juan } { mm } {
	\clearpage
	\tl_set:Nn \l__zhvt_chapter_title_tl { #1 }
	\tl_set:Nn \l_tmpa_tl { #2 }
	\tl_set_eq:NN \l_tmpc_tl \l_tmpa_tl
	\int_step_inline:nn { 7 - \str_count:n { #2 } }
	  { \tl_put_right:Nn \l_tmpc_tl { 　 } }
	\tl_set:Nx \l_tmpb_tl { \zhnum { page } }
	\tl_put_left:Nn \l_tmpb_tl { \hspace { 4em } }
	\tl_put_left:Nn \l_tmpb_tl { \l_tmpc_tl }
	\tl_put_left:Nn \l_tmpb_tl { \par \hspace { 3.5em } }
	\addtocontents { toc } { \l_tmpb_tl }
	\skip_horizontal:n { \l__zhvt_font_size_dim * 4 }
	\l_tmpa_tl
	\par
	\thispagestyle { empty }
}
\cs_set_protected:Npn \__zhvt_set_shukou_content: {
	% ★ 鱼尾上方留白（title 静态框）嵌入闲叟钤印：
	%   整页 /Rotate 90，故此处把图用 \rotatebox{90} 转正为竖章；
	%   尺寸卡在版心条幅内（书内约 7mm 宽 × 33.8mm 高，留余量不压鱼尾）。
	\setstaticcontents * { title } {
		\vbox_to_ht:nn { \l__zhvt_baseline_skip_dim } {
			\vskip 14.2pt
			\hbox_to_wd:nn { \l__zhvt_shukou_title_height_dim } {
				\hfil
				\rotatebox [ origin = c ] { 90 } { \includegraphics [ width = 7mm , height = 33.8mm ] { yshs.png } }
				\hfil
			}
		}
	}
% ★ 版心（书口）统一改法（2026-11-16 愚大人定）：18pt＋框内居中＋书口文字同版框色
%   ① 原 27.2pt 时 7 字长篇名（如「廉颇蔺相如列传」）末字溢出 chapter 框、压住页码；
%      改 \fontsize{18pt}{18pt}，前后 \hfil 令字串在框内居中（不随篇名字数漂移）。
%   ② \vbox{\vskip 8.2pt} 定量下移，抵消缩字号后栏宽方向 8.2pt 上偏（像素实测 y 中心 543.3 = 原版）。
%      ⚠ \raisebox / \vspace* 在 flowfram 框内无效；换字号须同步换补偿：16pt→10.1、19pt→7.4。
%   ③ chapter 与 page~number 均置 \color{sealbrown}（HTML 874434），与版框／界行同色；正文仍黑。
%   ⚠ 页码字位、钤印、鱼尾、界行、正文一律未动；五部合订本同此口径。
	\setdynamiccontents * { chapter }
	{ \color{sealbrown} \vbox { \vskip 8.2pt \hbox to \l__zhvt_shukou_chapter_title_height_dim
		{ \hfil \fontsize{18pt}{18pt}\selectfont \l__zhvt_title_str \quad \l__zhvt_chapter_title_tl \hfil } } }
	\setdynamiccontents * { page~number } { \color{sealbrown} \quad \zhnum { page } }
}
\__zhvt_set_shukou_content:
\ExplSyntaxOff

\maketitle

% 单篇，无多卷目录，不渲染目录页

\juan{卫康叔世家}{卷三十七~卫康叔世家第七}
\hspace*{2em}卫康叔名封。周武王同母少弟也。其次尚有冉季。冉季最少。\\
\hspace*{2em}武王已克殷纣。复以殷余民封纣子武庚禄父。比诸侯。以奉其先祀勿绝。为武庚未集。恐其有贼心。武王乃令其弟管叔。蔡叔傅相武庚禄父。以和其民。武王既崩。成王少。周公旦代成王治。当国。管叔。蔡叔疑周公。乃与武庚禄父作乱。欲攻成周。周公旦以成王命兴师伐殷。杀武庚禄父。管叔。放蔡叔。以武庚殷余民封康叔为卫君。居河。淇间故商墟。\\
\hspace*{2em}周公旦惧康叔齿少。乃申告康叔曰。﹁必求殷之贤人君子长者。问其先殷所以兴。所以亡。而务爱民。﹂告以纣所以亡者以淫于酒。酒之失。妇人是用。故纣之乱自此始。为梓材。示君子可法则。故谓之康诰酒诰梓材以命之。康叔之国。既以此命。能和集其民。民大说。\\
\hspace*{2em}成王长。用事。举康叔为周司寇。赐卫宝祭器。以章有德。\\
\hspace*{2em}康叔卒。子康伯代立。康伯卒。子考伯立。考伯卒。子嗣伯立。嗣伯卒。子伯立。伯卒。子靖伯立。靖伯卒。子贞伯立。贞伯卒。子顷侯立。\\
\hspace*{2em}顷侯厚赂周夷王。夷王命卫为侯。顷侯立十二年卒。子厘侯立。\\
\hspace*{2em}厘侯十三年。周厉王出奔于彘。共和行政焉。二十八年。周宣王立。\\
\hspace*{2em}四十二年。厘侯卒。太子共伯余立为君。共伯弟和有宠于厘侯。多予之赂。和以其赂赂士。以袭攻共伯于墓上。共伯入厘侯羡自杀。卫人因葬之厘侯旁。谥曰共伯。而立和为卫侯。是为武公。\\
\hspace*{2em}武公即位。修康叔之政。百姓和集。四十二年。犬戎杀周幽王。武公将兵往佐周平戎。甚有功。周平王命武公为公。五十五年。卒。子庄公扬立。\\
\hspace*{2em}庄公五年。取齐女为夫人。好而无子。又取陈女为夫人。生子。蚤死。陈女女弟亦幸于庄公。而生子完。完母死。庄公令夫人齐女子之。立为太子。庄公有宠妾。生子州吁。十八年。州吁长。好兵。庄公使将。石碏谏庄公曰。﹁庶子好兵。使将。乱自此起。﹂不听。二十三年。庄公卒。太子完立。是为桓公。\\
\hspace*{2em}桓公二年。弟州吁骄奢。桓公绌之。州吁出奔。十三年。郑伯弟段攻其兄。不胜。亡。而州吁求与之友。十六年。州吁收聚卫亡人以袭杀桓公。州吁自立为卫君。为郑伯弟段欲伐郑。请宋。陈。蔡与俱。三国皆许州吁。州吁新立。好兵。弑桓公。卫人皆不爱。石碏乃因桓公母家于陈。详为善州吁。至郑郊。石碏与陈侯共谋。使右宰丑进食。因杀州吁于濮。而迎桓公弟晋于邢而立之。是为宣公。\\
\hspace*{2em}宣公七年。鲁弑其君隐公。九年。宋督弑其君殇公。及孔父。十年。晋曲沃庄伯弑其君哀侯。\\
\hspace*{2em}十八年。初。宣公爱夫人夷姜。夷姜生子伋。以为太子。而令右公子傅之。右公子为太子取齐女。未入室。而宣公见所欲为太子妇者好。说而自取之。更为太子取他女。宣公得齐女。生子寿。子朔。令左公子傅之。太子伋母死。宣公正夫人与朔共谗恶太子伋。宣公自以其夺太子妻也。心恶太子。欲废之。及闻其恶。大怒。乃使太子伋于齐而令盗遮界上杀之。与太子白旄。而告界盗见持白旄者杀之。且行。子朔之兄寿。太子异母弟也。知朔之恶太子而君欲杀之。乃谓太子曰。﹁界盗见太子白旄。即杀太子。太子可毋行。﹂太子曰。﹁逆父命求生。不可。﹂遂行。寿见太子不止。乃盗其白旄而先驰至界。界盗见其验。即杀之。寿已死。而太子伋又至。谓盗曰。﹁所当杀乃我也。﹂盗并杀太子伋。以报宣公。宣公乃以子朔为太子。十九年。宣公卒。太子朔立。是为惠公。\\
\hspace*{2em}左右公子不平朔之立也。惠公四年。左右公子怨惠公之谗杀前太子伋而代立。乃作乱。攻惠公。立太子伋之弟黔牟为君。惠公奔齐。\\
\hspace*{2em}卫君黔牟立八年。齐襄公率诸侯奉王命共伐卫。纳卫惠公。诛左右公子。卫君黔牟奔于周。惠公复立。惠公立三年出亡。亡八年复入。与前通年凡十三年矣。\\
\hspace*{2em}二十五年。惠公怨周之容舍黔牟。与燕伐周。周惠王奔温。卫。燕立惠王弟颓为王。二十九年。郑复纳惠王。三十一年。惠公卒。子懿公赤立。\\
\hspace*{2em}懿公即位。好鹤。淫乐奢侈。九年。翟伐卫。卫懿公欲发兵。兵或畔。大臣言曰。﹁君好鹤。鹤可令击翟。﹂翟于是遂入。杀懿公。\\
\hspace*{2em}懿公之立也。百姓大臣皆不服。自懿公父惠公朔之谗杀太子伋代立至于懿公。常欲败之。卒灭惠公之后而更立黔牟之弟昭伯顽之子申为君。是为戴公。\\
\hspace*{2em}戴公申元年卒。齐桓公以卫数乱。乃率诸侯伐翟。为卫筑楚丘。立戴公弟燬为卫君。是为文公。文公以乱故奔齐。齐人入之。\\
\hspace*{2em}初。翟杀懿公也。卫人怜之。思复立宣公前死太子伋之后。伋子又死。而代伋死者子寿又无子。太子伋同母弟二人。其一曰黔牟。黔牟尝代惠公为君。八年复去。其二曰昭伯。昭伯。黔牟皆已前死。故立昭伯子申为戴公。戴公卒。复立其弟燬为文公。\\
\hspace*{2em}文公初立。轻赋平罪。身自劳。与百姓同苦。以收卫民。\\
\hspace*{2em}十六年。晋公子重耳过。无礼。十七年。齐桓公卒。二十五年。文公卒。子成公郑立。\\
\hspace*{2em}成公三年。晋欲假道于卫救宋。成公不许。晋更从南河度。救宋。征师于卫。卫大夫欲许。成公不肯。大夫元咺攻成公。成公出奔。晋文公重耳伐卫。分其地予宋。讨前过无礼及不救宋患也。卫成公遂出奔陈。二岁。如周求入。与晋文公会。晋使人鸩卫成公。成公私于周主鸩。令薄。得不死。已而周为请晋文公。卒入之卫。而诛元咺。卫君瑕出奔。七年。晋文公卒。十二年。成公朝晋襄公。十四年。秦穆公卒。二十六年。齐邴歜弑其君懿公。三十五年。成公卒。子穆公遬立。\\
\hspace*{2em}穆公二年。楚庄王伐陈。杀夏征舒。三年。楚庄王围郑。郑降。复释之。十一年。孙良夫救鲁伐齐。复得侵地。穆公卒。子定公臧立。定公十二年卒。子献公衎立。\\
\hspace*{2em}献公十三年。公令师曹教宫妾鼓琴。妾不善。曹笞之。妾以幸恶曹于公。公亦笞曹三百。十八年。献公戒孙文子。宁惠子食。皆往。日旰不召。而去射鸿于囿。二子从之。公不释射服与之言。二子怒。如宿。孙文子子数侍公饮。使师曹歌巧言之卒章。师曹又怒公之尝笞三百。乃歌之。欲以怒孙文子。报卫献公。文子语蘧伯玉。伯玉曰。﹁臣不知也。﹂遂攻出献公。献公奔齐。齐置卫献公于聚邑。孙文子。宁惠子共立定公弟秋为卫君。是为殇公。\\
\hspace*{2em}殇公秋立。封孙文子林父于宿。十二年。宁喜与孙林父争宠相恶。殇公使宁喜攻孙林父。林父奔晋。复求入故卫献公。献公在齐。齐景公闻之。与卫献公如晋求入。晋为伐卫。诱与盟。卫殇公会晋平公。平公执殇公与宁喜而复入卫献公。献公亡在外十二年而入。\\
\hspace*{2em}献公后元年。诛宁喜。\\
\hspace*{2em}三年。吴延陵季子使过卫。见蘧伯玉。史鳅。曰。﹁卫多君子。其国无故。﹂过宿。孙林父为击磬。曰。﹁不乐。音大悲。使卫乱乃此矣。﹂是年。献公卒。子襄公恶立。\\
\hspace*{2em}襄公六年。楚灵王会诸侯。襄公称病不往。\\
\hspace*{2em}九年。襄公卒。初。襄公有贱妾。幸之。有身。梦有人谓曰。﹁我康叔也。令若子必有卫。名而子曰﹃元﹄。﹂妾怪之。问孔成子。成子曰。﹁康叔者。卫祖也。﹂及生子。男也。以告襄公。襄公曰。﹁天所置也。﹂名之曰元。襄公夫人无子。于是乃立元为嗣。是为灵公。\\
\hspace*{2em}灵公五年。朝晋昭公。六年。楚公子弃疾弑灵王自立。为平王。十一年。火。\\
\hspace*{2em}三十八年。孔子来。禄之如鲁。后有隙。孔子去。后复来。\\
\hspace*{2em}三十九年。太子蒯聩与灵公夫人南子有恶。欲杀南子。蒯聩与其徒戏阳遬谋。朝。使杀夫人。戏阳后悔。不果。蒯聩数目之。夫人觉之。惧。呼曰。﹁太子欲杀我。﹂灵公怒。太子蒯聩奔宋。已而之晋赵氏。\\
\hspace*{2em}四十二年春。灵公游于郊。令子郢仆。郢。灵公少子也。字子南。灵公怨太子出奔。谓郢曰。﹁我将立若为后。﹂郢对曰。﹁郢不足以辱社稷。君更图之。﹂夏。灵公卒。夫人命子郢为太子。曰。﹁此灵公命也。﹂郢曰。﹁亡人太子蒯聩之子辄在也。不敢当。﹂于是卫乃以辄为君。是为出公。\\
\hspace*{2em}六月乙酉。赵简子欲入蒯聩。乃令阳虎许命卫十余人衰绖归。简子送蒯聩。卫人闻之。发兵击蒯聩。蒯聩不得入。入宿而保。卫人亦罢兵。\\
\hspace*{2em}出公辄四年。齐田乞弑其君孺子。八年。齐鲍子弑其君悼公。\\
\hspace*{2em}孔子自陈入卫。九年。孔文子问兵于仲尼。仲尼不对。其后鲁迎仲尼。仲尼反鲁。\\
\hspace*{2em}十二年。初。孔圉文子取太子蒯聩之姊。生悝。孔氏之竖浑良夫美好。孔文子卒。良夫通于悝母。太子在宿。悝母使良夫于太子。太子与良夫言曰。﹁苟能入我国。报子以乘轩。免子三死。毋所与。﹂与之盟。许以悝母为妻。闰月。良夫与太子入。舍孔氏之外圃。昏。二人蒙衣而乘。宦者罗御。如孔氏。孔氏之老栾宁问之。称姻妾以告。遂入。适伯姬氏。既食。悝母杖戈而先。太子与五人介。舆豭从之。伯姬劫悝于厕。强盟之。遂劫以登台。栾宁将饮酒。炙未熟。闻乱。使告仲由。召护驾乘车。行爵食炙。奉出公辄奔鲁。\\
\hspace*{2em}仲由将入。遇子羔将出。曰。﹁门已闭矣。﹂子路曰。﹁吾姑至矣。﹂子羔曰。﹁不及。莫践其难。﹂子路曰。﹁食焉不辟其难。﹂子羔遂出。子路入。及门。公孙敢阖门。曰。﹁毋入为也。﹂子路曰。﹁是公孙也。求利而逃其难。由不然。利其禄。必救其患。﹂有使者出。子路乃得入。曰。﹁太子焉用孔悝。虽杀之。必或继之。﹂且曰。﹁太子无勇。若燔台。必舍孔叔。﹂太子闻之。惧。下石乞。盂黡敌子路。以戈击之。割缨。子路曰。﹁君子死。冠不免。﹂结缨而死。孔子闻卫乱。曰。﹁嗟乎。柴也其来乎。由也其死矣。﹂孔悝竟立太子蒯聩。是为庄公。\\
\hspace*{2em}庄公蒯聩者。出公父也。居外。怨大夫莫迎立。元年即位。欲尽诛大臣。曰。﹁寡人居外久矣。子亦尝闻之乎。﹂群臣欲作乱。乃止。\\
\hspace*{2em}二年。鲁孔丘卒。\\
\hspace*{2em}三年。庄公上城。见戎州。曰。﹁戎虏何为是。﹂戎州病之。十月。戎州告赵简子。简子围卫。十一月。庄公出奔。卫人立公子斑师为卫君。齐伐卫。虏斑师。更立公子起为卫君。\\
\hspace*{2em}卫君起元年。卫石曼尃逐其君起。起奔齐。卫出公辄自齐复归立。初。出公立十二年亡。亡在外四年复入。出公后元年。赏从亡者。立二十一年卒。出公季父黔攻出公子而自立。是为悼公。\\
\hspace*{2em}悼公五年卒。子敬公弗立。敬公十九年卒。子昭公纠立。是时三晋强。卫如小侯。属之。\\
\hspace*{2em}昭公六年。公子亹弑之代立。是为怀公。怀公十一年。公子颓弑怀公而代立。是为慎公。慎公父。公子适。适父。敬公也。慎公四十二年卒。子声公训立。声公十一年卒。子成侯遬立。\\
\hspace*{2em}成侯十一年。公孙鞅入秦。十六年。卫更贬号曰侯。\\
\hspace*{2em}二十九年。成侯卒。子平侯立。平侯八年卒。子嗣君立。\\
\hspace*{2em}嗣君五年。更贬号曰君。独有濮阳。\\
\hspace*{2em}四十二年卒。子怀君立。怀君三十一年。朝魏。魏囚杀怀君。魏更立嗣君弟。是为元君。元君为魏婿。故魏立之。元君十四年。秦拔魏东地。秦初置东郡。更徙卫野王县。而并濮阳为东郡。二十五年。元君卒。子君角立。\\
\hspace*{2em}君角九年。秦并天下。立为始皇帝。二十一年。二世废君角为庶人。卫绝祀。\\
\hspace*{2em}太史公曰。余读世家言。至于宣公之太子以妇见诛。弟寿争死以相让。此与晋太子申生不敢明骊姬之过同。俱恶伤父之志。然卒死亡。何其悲也。或父子相杀。兄弟相灭。亦独何哉。\\

% ===== 索隐述赞（司马贞《史记索隐》）· 统一样式 =====
% 题头「索隐述赞曰」五字并入反白框；正文每四句一列，定宽盒 \hfil 居列中；
% \fontsize{24}{41} 列距对齐全书界行（121+41k）；\clearpage 使述赞独占一页。
\clearpage
\shuzanhead{索隐述赞曰}\linebreak
{\shuzanbody\fontsize{24}{41}\selectfont
\hbox to \hsize{\hfil 司寇受封。梓材有作。成锡厥器。夷加其爵。\hfil}\linebreak
\hbox to \hsize{\hfil 暨武能脩。从文始约。诗美归燕。传矜石碏。\hfil}\linebreak
\hbox to \hsize{\hfil 皮冠射鸿。乘轩使穀。宣纵淫嬖。衅生伋。\hfil}\linebreak
\hbox to \hsize{\hfil 朔。蒯聩得罪。出公行恶。卫祚日衰。\hfil}\linebreak
\hbox to \hsize{\hfil 失於君角。\hfil}%
\par}
\end{document}
